\documentclass[12pt, spanish]{article}
\usepackage[spanish]{babel}
\selectlanguage{spanish}
%\usepackage{natbib}
\usepackage{url}
\usepackage[utf8x]{inputenc}
\usepackage{graphicx}
\graphicspath{{images/}}
\usepackage{parskip}
\usepackage{fancyhdr}
\usepackage{vmargin}
\usepackage{multirow}
\usepackage{float}
\usepackage{chngpage}
\usepackage{enumitem}
\usepackage{forloop}


\usepackage{amsfonts}
\usepackage{amsmath}

\usepackage{subcaption}

\usepackage{hyperref}
\usepackage[
    type={CC},
    modifier={by-nc-sa},
    version={4.0},
]{doclicense}

\hypersetup{
    colorlinks=true,
    linkcolor=blue,
    filecolor=magenta,
    urlcolor=cyan,
}

% para codigo
\usepackage{listings}
\usepackage{xcolor}


\usepackage[default]{sourcesanspro}

\setmarginsrb{2 cm}{1 cm}{2 cm}{2 cm}{1 cm}{1.5 cm}{1 cm}{1.5 cm}

\title{Práctica 1:\\
Aritmética modular.\hspace{0.05cm} }
\author{Antonio David Villegas Yeguas}
\date{\today}

\makeatletter
\let\thetitle\@title
\let\theauthor\@author
\let\thedate\@date
\makeatother

\pagestyle{fancy}
\fancyhf{}
\rhead{\theauthor}
\lhead{\thetitle}
\cfoot{\thepage}


\usepackage{caption}

\lstset{
language=C++,
basicstyle=\small\ttfamily,
numbers=left,
numbersep=5pt,
xleftmargin=20pt,
frame=tb,
framexleftmargin=20pt
}


\DeclareCaptionFormat{mylst}{\hrule#2#3}
\captionsetup[lstlisting]{format=mylst,labelfont=bf,singlelinecheck=off,labelsep=space}

\begin{document}

%%%%%%%%%%%%%%%%%%%%%%%%%%%%%%%%%%%%%%%%%%%%%%%%%%%%%%%%%%%%%%%%%%%%%%%%%%%%%%%%%%%%%%%%%

\begin{titlepage}
    \centering
    \vspace*{0.3 cm}
    \includegraphics[scale = 0.50]{ugr.png}\\[0.7 cm]
    %\textsc{\LARGE Universidad de Granada}\\[2.0 cm]
    \textsc{\large 4º CSI 2020/21}\\[0.5 cm]
    \textsc{\large Grado en Ingeniería Informática}\\[0.5 cm]
    \rule{\linewidth}{0.2 mm} \\[0.2 cm]
    { \huge \bfseries \thetitle}\\
    \rule{\linewidth}{0.2 mm} \\[1 cm]

    \begin{minipage}{0.4\textwidth}
        \begin{flushleft} \large
            \emph{Autor:}\\
            \theauthor\\
			 \emph{DNI:}\\
            77021623-M
            \end{flushleft}
            \end{minipage}~
            \begin{minipage}{0.4\textwidth}
            \begin{flushright} \large
            \emph{Asignatura: \\
            Criptografía}   \\
            \emph{Correo:}\\
            advy99@correo.ugr.es
        \end{flushright}
    \end{minipage}\\[0.5cm]

    {\large \thedate}\\[0.5cm]
    {\url{https://github.com/advy99/CRIP/}}
    {\doclicenseThis}

    \vfill

\end{titlepage}

%%%%%%%%%%%%%%%%%%%%%%%%%%%%%%%%%%%%%%%%%%%%%%%%%%%%%%%%%%%%%%%%%%%%%%%%%%%%%%%%%%%%%%%%%

\tableofcontents
\pagebreak

%%%%%%%%%%%%%%%%%%%%%%%%%%%%%%%%%%%%%%%%%%%%%%%%%%%%%%%%%%%%%%%%%%%%%%%%%%%%%%%%%%%%%%%%%


\section*{Introducción}

En esta práctica trabajaremos con funciones de un solo sentido de cara a comprender su uso para cifrar y firmar mensajes y hacer un resumen de mensajes, así como implementar algunas de las más utilizadas como la función de Rabin o RSA.


\section{Cifrado utilizando la función knapsack}

Este método consiste en generar una secuencia super-creciente (la suma de los terminos que preceden a un elemento es menor que dicho elemento). Con esta secuencia, para cifrar un mensaje en binario simplemente se escogen los elementos de la secuencia super-creciente en los que la posición asociada en el mensaje esté como verdadero.

Esto tiene la ventaja de que quien conozca la secuencia super-creciente puede resolver este problema utilizando un algoritmo voraz en muy pocos pasos, la longitud de la secuencia, sin embargo, si la secuencia no es supre-creciente este problema se trata de un problema NP completo, siendo imposible de resolver de forma práctica por un algoritmo conocido.

Sabiendo esto, el objetivo será que la secuencia super-creciente forme parte de la clave privada, de forma que quien la conozca pueda obtener los mensajes de forma rápida, y se utilice como clave pública una transformación a la secuencia super-creciente, de forma que escogeremos un $n$ mayor que la suma de la secuencia, y una transformación en $Z_n$ haga que la secuencia no sea super-creciente, convirtiéndolo en un problema NP completo.

Esto lo haremos escogiendo un $u$ tal que $gcd(n,u) = 1$ de cara a que $u$ tenga inverso en $Z_n$. Con esto, para un elemento de la secuencia super-creciente $a_i$, definiremos $a_i^* = u \cdot a_i mod n$, de forma que la secuencia $(a_1^*, \dots, a_k^*)$ será la clave pública y $((a_1, \dots, a_k),n, u)$ será la clave privada.

De esta forma, si alguien quiere mandarnos un mensaje solo tiene que cifrarlo con la clave pública:

$$
\sum_{i = 1}^{k} x_i a_i^*
$$


\begin{lstlisting}[caption={Algoritmo para descifrar un mensaje}]
INPUT: secuencia de bits (mensaje): m, clave puvlica: clave
OUTPUT: secuencia de bits que conforman el mensaje

resultado <- 0

Para cada elemento de m:
	resultado <- resultado + ( m[i] * clave[i])

Devolver resultado
\end{lstlisting}


Y para descifrar será necesario calcular el inverso de $u$ en $Z_n$, al número dado (mensaje cifrado), lo multiplicamos por el inverso de $u$, de forma que este número que se había calculado multiplicando por $u$, al multiplicar por su inverso obtenemos el número como si lo hubiéramos obtenido con la suma de elementos de la secuencia super-creciente, y no la secuencia transformada a $Z_n$.

De esta forma, como conocemos la secuencia super-creciente podemos obtener con un algoritmo voraz el mensaje original:

\begin{lstlisting}[caption={Algoritmo voraz para la función knapsack}]
INPUT: secuencia super-creciente: s, numero a descifrar: numero
OUTPUT: secuencia de bits que conforman el mensaje

resultado = [0,...,0] <- misma longitud que s

suma <- 0
i <- 0
Mientras suma != numero:
	Si suma + s[i] <= numero:
		resultado[i] <- 1
		suma <- suma + numero
	i <- i+1

Devolver resultado
\end{lstlisting}

La función para descifrar quedaría de la siguiente forma:

\begin{lstlisting}[caption={Algoritmo para descifrar un mensaje}]
INPUT: secuencia super-creciente: s, numero a descifrar: numero, n, u
OUTPUT: secuencia de bits que conforman el mensaje

inverso_u <- inverso(u, n) // inverso de u mod n
numero_a_descifrar <- (numero * inverso_u) mod n

Devolver algoritmo_voraz_knapsack(s, numero_a_descifrar)
\end{lstlisting}

Algunos ejemplos de la ejecución:

\begin{lstlisting}
➜ ./bin/ejercicio1 13 1010100100011
Llave publica:
92627 32072 68013 67453 90667 56075 23392 2545 61496 26216 44414 16106 24194
Llave privada:
17 20 56 110 206 429 857 1700 3408 6810 13615 27239 54473
La llave privada generada es una secuencia super-creciente

N: 108943
U: 88758
Mensaje cifrado: 294152
Mensaje descifrado: 1010100100011
Mensaje       dado: 1010100100011


➜ ./bin/ejercicio1 20 10100010001001010110
Llave publica:
12475022 4832564 10742065 12633531 5149582 8828615 6259145 6346381 6258445 13856235 12558758 2190142 3172143 11439258 5046114 1110425 2220850 4441700 8883400 3821229
Llave privada:
17 20 56 110 206 429 857 1700 3408 6810 13615 27239 54473 108944 217895 435799 871598 1743196 3486392 6972784
La llave privada generada es una secuencia super-creciente

N: 13945571
U: 11398085
Mensaje cifrado: 67909773
Mensaje descifrado: 10100010001001010110
Mensaje       dado: 10100010001001010110
\end{lstlisting}

Vemos como la clave privada es facil reconocer que es super-creciente, mientras que la pública parecen una secuencia de números sin correlación.

\section{Obteniendo y en ElGammal conociendo p, alfa y la imagen de x}

En este ejercicio se nos pide escoger un primo $p$ mayor que nuestro DNI, encontrar un elemento primitivo $\alpha$ en $Z_p$. Con esto, se nos da la función para obtener el valor $y$ de la clave pública de ElGammal, además de pedirnos usar nuestra fecha de nacimiento como número a calcular el inverso de la función.

Para resolver este ejercicio he escogido un primo $p$ que sea mayor que un número dado, y que $\frac{p}{2} - 1$ también sea primo. De esta forma, para calcular $\alpha$, simplemente he buscado un valor que el símbolo de Jacobi de dicho valor y $p$ sea igual a -1.

Con esto, sabiendo el valor de $\alpha$, el valor de la imagen y $p$, he utilizado el logaritmo discreto para obtener el inverso del número pedido:

\begin{lstlisting}[caption={Ejercicio 2}]
INPUT: numero para calcular p: posible_p, numero a calcular el inverso: a_invertir
OUTPUT: inverso de a_invertir

p <- siguiente_primo(posible_p)

Mientras (p / 2) - 1 no sea primo:
	p <- siguiente_primo(p)

alpha <- 2

Mientras simbolo_jacobi(alpha, p) != -1:
	alpha <- alpha + 1

x <- logaritmo_discreto(alpha, a_invertir, p)

Devolver x
\end{lstlisting}

La ejecución con mi implementación:

\begin{lstlisting}
➜ ./bin/ejercicio2 77021623 19990915
Utilizaré el primo: 77022299
Un elemento primitivo de Z_77022299 es 2
El x tal que 2^x en Z_77022299 = 19990915 es 28098736
Comprobación:
Fecha introducida: 19990915
Valor calculado con alpha ^ x: 19990915
\end{lstlisting}


\section{Obtener p y q dado un n y dos valores con la misma imagen para la función de Rabin}

En este ejercicio se nos da la función de Rabin: $f(x) = x^2$, también se nos da $n = 48478872564493742276963$, y que $f (12) =144 = f (37659670402359614687722)$, y con esto tenemos que calcular los factores primos $p$ y $q$ de $n$.

Para resolver este ejercicio nos basaremos en que hemos encontrado dos números $x$ e $y$ con la misma función de Rabin, pero que no son iguales. Con esto sabemos que $12^2 \cong y^2 mod n$, por lo que $(x + y)(x - y) = 0$, por lo que $n$ divide a $(x + y)(x - y)$, pero no a $(x - y)$, ya que $x$ e $y$ son distintos, luego $gcd(n, x - y)$ es un factor no trivial.

Para este ejercicio simplemente calculamos el máximo común divisor de $n$ con $x - y$ y de $n$ con $x - (n - y)$ para obtener $p$ y $q$:

\begin{lstlisting}
➜ ./bin/ejercicio3 48478872564493742276963 12 37659670402359614687722
p = 159497098847
q = 303948303229
p * q         = 48478872564493742276963
n introducido = 48478872564493742276963
\end{lstlisting}


\section{Función resumen utilizando la construcción de Merkle-Damgar}

\section{Descifrar RSA conociendo p y q}

\section{Descifrar RSA conociendo la clave publica y privada}

\section{Implementando la firma RSA}

\subsection{Generación de llaves}

\subsection{Generación de la firma asociada a un fichero}

\subsection{Verificación de la firma asociada a un fichero}






\end{document}
